\section{Две леммы и вычисляемый шаг}
Две леммы этого раздела используют \eqref{eq:gap_smoothness},
но не требуют выпуклости или квазар-выпуклости функции $f$.

\begin{lemma}[Рост зазора и верхняя модель]\label{lem:transport}
Пусть $s,d$ заданы в \eqref{eq:gap}. Для $0\le t\le1$ верно
\begin{align}
 G(x+td)&\le(1-t)G(x)
    +D\|d\|\int_0^t\bigl(L_0+L_1G(x+ud)\bigr)\,du,
       \label{eq:transport}\\
 L_0+L_1G(x+td)&\le
      \bigl(L_0+L_1G(x)\bigr)e^{L_1D\|d\|t}.
       \label{eq:transport_exp}
\end{align}
Поэтому для $0\le\alpha\le1$
\begin{equation}\label{eq:upper_model}
 f(x+\alpha d)\le f(x)-\alpha G(x)
 +(L_0+L_1G(x))\|d\|^2
       \int_0^\alpha(\alpha-t)e^{L_1D\|d\|t}\,dt.
\end{equation}
\end{lemma}
\begin{proof}
Для любого $v\in Q$ раскроем скалярное произведение относительно $x$:
\[
 \langle\nabla f(x+td),x+td-v\rangle
 \le G(x)-tG(x)
       +D\|\nabla f(x+td)-\nabla f(x)\|.
\]
Берём максимум по $v$ и интегрируем гессиан вдоль отрезка.
Получаем \eqref{eq:transport}. Отбрасываем неположительное
слагаемое $-tG(x)$ и применяем неравенство Гронуолла к
$L_0+L_1G(x+td)$; это даёт \eqref{eq:transport_exp}.
При $L_1=0$ оценка непосредственна.
Формула Тейлора с интегральным остатком даёт \eqref{eq:upper_model},
поскольку
$|d^\top\nabla^2f(x+td)d|\le\|\nabla^2f(x+td)\|\|d\|^2$.
\end{proof}

Экспонента в \eqref{eq:upper_model} ограничивает рост кривизны
при движении. Квадрат $\|d\|^2$ возникает из второй производной
вдоль направления. Оба множителя получены в доказательстве.
Выбранное направление существенно. Например, для
$f(t)=t+t^3/6$ на $[0,1]$ выполнено $|f''(t)|\le G(t)$,
однако $G(0)=0$ и $f(1)>f(0)+f'(0)$.
Поэтому верхнюю модель с начальным коэффициентом $L_0+L_1G(x)$
нельзя без изменения переносить на произвольную пару точек $x,y$.

\paragraph{Правило шага.}
При $G(x)>0$ и $L_1>0$ минимизируем правую часть
\eqref{eq:upper_model} на $[0,1]$:
\begin{equation}\label{eq:log_step}
 \boxed{\displaystyle
 \alpha=\min\left\{1,\,
 \frac{1}{L_1D\|d\|}
 \log\left(1+
 \frac{L_1D\,G(x)}
 {(L_0+L_1G(x))\|d\|}\right)\right\}.}
\end{equation}
Здесь и далее $\log$ означает натуральный логарифм.
При $L_1=0<L_0$ используем непрерывный предел
$\alpha=\min\{1,G(x)/(L_0\|d\|^2)\}$.
При $L_0=L_1=0$ берём $\alpha=1$: функция на $Q$ аффинна.
При $G(x)=0$ останавливаемся до вычисления дробей.

\begin{lemma}[Спуск]\label{lem:descent}
Указанный шаг удовлетворяет
\begin{equation}\label{eq:descent}
 f(x+\alpha d)\le f(x)-\frac{\alpha G(x)}2.
\end{equation}
Кроме отдельно разобранного аффинного случая,
\begin{equation}\label{eq:step_lower}
 \alpha\ge
 \min\left\{1,
   \frac{(\log2)G(x)}{(L_0+L_1G(x))D^2}\right\}.
\end{equation}
\end{lemma}
\begin{proof}
Только для этого скалярного вычисления положим
$b=L_1D\|d\|>0$. Интеграл в \eqref{eq:upper_model} равен
$(e^{b\alpha}-1-b\alpha)/b^2$.
Производная верхней модели имеет вид
\[
 -G(x)+(L_0+L_1G(x))\|d\|^2
             \frac{e^{b\alpha}-1}{b}.
\]
Её ноль даёт \eqref{eq:log_step}. В минимизаторе с ограничением
$\alpha\le1$ производная неположительна.
Функция $(e^{bt}-1)/b$ выпукла и равна нулю при $t=0$.
Поэтому её интеграл на $[0,\alpha]$ не больше половины произведения
$\alpha$ на значение в правом конце. Подставляем это в
\eqref{eq:upper_model} и получаем \eqref{eq:descent}.

Для \eqref{eq:step_lower} временно положим
$r=\|d\|/D\le1$ и $q=L_1G(x)/(L_0+L_1G(x))\le1$.
Неограниченный шаг равен $\log(1+q/r)/(L_1D^2r)$, причём
\[
 \log(1+q/r)\ge\log(1+q)\ge q\log2\ge qr\log2.
\]
Среднее неравенство следует из вогнутости логарифма.
При $L_1=0$ оба утверждения следуют из квадратичной модели.
\end{proof}

В шаге через норму градиента из~\cite{vyguzov} длина шага
ограничивает экспоненту числом $e$. Формула~\eqref{eq:log_step}
сразу минимизирует экспоненциальную модель. Логарифм обеспечивает
\[
 e^{L_1D\|d\|\alpha}-1
 \le\frac{L_1D\,G(x)}{(L_0+L_1G(x))\|d\|}.
\]
Именно это неравенство используется в доказательстве спуска.

\paragraph{Алгоритм и остановка.}
Зададим точность $\varepsilon>0$.
В каждой посещённой точке $x_k$ вычисляем $f(x_k)$, градиент
и линейный минимум по всему $Q$. Сохраняем
\begin{equation}\label{eq:stop}
 \ell_k=\max_{0\le j\le k}
        \left\{f(x_j)-\frac{G(x_j)}{\gamma}\right\},
 \qquad \text{остановка при } f(x_k)-\ell_k\le\varepsilon.
\end{equation}
Иначе полагаем $x_{k+1}=x_k+\alpha_kd_k$, используя
\eqref{eq:log_step} или указанные предельные случаи.
Каждое новое значение функции вычисляется один раз. Подбора шага нет.
